%% kapa: exam=1 round=22 year=2011 period=2 form=A subject=accounting subject_ko=회계학 seq=18 number=18
%% source: Q-Net 감정평가사 기출문제 (cst00309) artlSeq=5036790 file=extracted/1008/5036790/2011년 제22회 감정평가사 1차시험 기출문제/2교시.hwp
%% conversion: kordoc 4.24.51 -> tools/gen1.py (auto); review=proofread-gpt-astra; flags=none
\begin{question}{18}
회계변경에 관한 설명으로 옳지 않은 것은?
\begin{choices}
  \item 기업이 하나의 일반적으로 인정된 회계원칙(GAAP)에서 다른 회계원칙(GAAP)으로 바꾸는 것을 회계정책의 변경이라 한다.
  \item 감가상각자산의 내용연수 또는 감가상각에 내재된 미래 경제적효익의 기대소비행태가 변하는 경우 회계정책의 변경으로 처리한다.
  \item 회계정책의 변경을 반영한 재무제표가 특정 거래, 기타 사건 또는 상황이 재무상태, 경영성과 또는 현금흐름에 미치는 영향에 대해서 신뢰성 있고 더 목적적합한 정보를 제공하는 경우 회계정책의 변경이 가능하다.
  \item 회계정책의 변경에 대해서는 소급법을 적용하여 재무제표를 작성한다.
  \item 회계정책의 변경과 회계추정의 변경을 구분하는 것이 어려운 경우에는 이를 회계추정의 변경으로 본다.
\end{choices}
\end{question}
\answer{2}
